% Pursuing Stacks (cotes 134-2 to 134-8): archive cote 134-4, site page 23 (PDF page 24), typescript page 111.
% typescript. Draft compared: ps3.tex:1124-1197.
% AI-transcribed by Claude Opus 5.5 (claude-opus-5-5, adaptive thinking, effort high), Message Batches API, batch ps2 (msgbatch_01NVwr1Yuo8qJf5rykeGBXPK), 2026-11-12; not yet checked by eye.
\scanpage{134-4/23}
\mspage{111}

\margin{51.} First thing though I would like to do now,  is to elaborate ``from
scratch'' on the four basic homotopy notions and their interplay,  much in
the style of a ``fugue with variations'' I guess,  and without interference
of a pregiven notion W of weak equivalence  -  relationship with a W will
be examined only,  after the intrinsic homotopy notions and their interre-%
lations are well understood.

We start with a category M ,  without for the time being any specific
assumptions on M.  T\hand{h}e strongest we're going to introduce,  I guess,  is
existence of finite products,  and incidently maybe finite sums and fiber
products.  In the cases we have in mind,  M is a ``large'' category,  there-%
fore it does'nt seem timely here to introduce M\^{} and the embedding of M
into M\^{}.

T\hand{h}e most trivial implications between the four basic \inserted[typed]{homotopy} notions are
symbolized by the plain arrows in the diagram below,  the somewhat more
technic\hand{$\frown$}al ones by dotted arrows \hand{*}.  It is understood these notions cor-%
respond to a given ``homotopy structure'' on M,  symbolized by the letter h,
and which (in  the most favorable cases)  may be described \inserted{at will} in terms of any
one of the four notions.  I'll first describe separately each of these
basic notions,  and afterwards the relationsships \xout{expressed} symbolized by
the arrows in the diagram.  I recall that an \uline{interval} in M is just an
object  I ,  endowed with two subobjects $e_o,e_1$ which are final objects of
M,  or equivalently,  with two sections $\hand{$\delta$}_o$ ,$\hand{$\delta$}_1$ of I over a fixed final
object $e_{\hand{M}}=e_{\deleted{\unsure{M}}}$ of M.  I definitely want to forget entirely for the time being
about any condition of the type $e_o\hand{$\cap$}e_1 = \emptyset_M$  (initial object of M)  -
we may later refer to these as ``\uline{separated}'' intervals (namely the endpoints
$e_o,e_1$ are ``separated'').  We'll denote by Int(M) the set of all intervals
in M ,  by \uline{Int}(M) the corresponding category (the notion of a morphism of
\margin{\unsure{$\llcorner$}}intervals being the obvious one).  Now here's the organigram:

\margin{(D)}
\diagram{
\begin{tikzcd}[arrows=Rightarrow, column sep=huge, row sep=large]
\begin{array}{l}\text{1\hand{)} homotopy relation}\\ R_h\hand{$\subset$}\mathrm{Fl}(M)\mathrm{x}\mathrm{Fl}(M)\end{array}
  \ar[r, "\hand{}"] \ar[d] &
\begin{array}{l}\text{2\hand{)} homotopisms}\\ W_h\hand{$\subset$}\mathrm{Fl}(M)\end{array}
  \ar[l, bend right=15, dashed, swap, "\hand{1}"] \ar[d] \\
\begin{array}{l}\text{3\hand{)} homotopy intervals}\\ \hand{$\Sigma$}_h\hand{$\subset$}\mathrm{Int}(M)\end{array}
  \ar[u, bend left=30, dashed, "\hand{2}"] &
\begin{array}{l}\text{4\hand{)} contractible objects}\\ C_h\hand{$\subset$}\mathrm{Ob}(M)\end{array}
  \ar[l, dashed, swap, "\hand{3}"]
\end{tikzcd}
}
\margin{*}

\margin{23}
